% =====================================================================
% Stage Report — BEV-VAWA-Lite
% Compile: pdflatex stage_report.tex  (twice, for cross-refs)
% Fonts  : Times New Roman via mathptmx (serif) + helvet (sans fallback).
% Colour : all black; hyperref links use black underlines, not colour.
% =====================================================================
\documentclass[11pt,a4paper]{article}

\usepackage[margin=1in]{geometry}
\usepackage{mathptmx}              % Times New Roman (text + math)
\usepackage[T1]{fontenc}
\usepackage[utf8]{inputenc}
\usepackage{graphicx}
\usepackage{booktabs}
\usepackage{multirow}
\usepackage{amsmath,amssymb}
\usepackage{xcolor}
\usepackage{microtype}
\usepackage[numbers,sort&compress]{natbib}
\usepackage[hidelinks]{hyperref}
\usepackage{caption}
\captionsetup{font=small,labelfont=bf}

% Enforce all-black text, including links / citations / section headings.
\color{black}
\hypersetup{colorlinks=false,
            allbordercolors={1 1 1},
            pdfborder={0 0 0},
            citecolor=black,
            linkcolor=black,
            urlcolor=black}

\title{\textbf{Stage Report: A Lightweight Vision-Action / World-Action\\
Navigation Stack with Reactive Safety Augmentation\\
on the PIB-Nav Benchmark}}
\author{BEV-VAWA-Lite Project}
\date{April 2026}

\begin{document}
\maketitle

\begin{abstract}
We report interim results on \textbf{BEV-VAWA-Lite}, a lightweight
depth-only PointGoal navigation stack combining a reactive Vision-Action
(VA) branch with a short-horizon latent World-Action (WA) branch over a
shared bird's-eye-view (BEV) state. The system is trained entirely on a
consumer-grade Apple M4 laptop using a procedural MuJoCo benchmark
(\emph{PIB-Nav}). This report consolidates four milestones: (i) a
full-scale retrain that raised Success Rate (SR) at our headline seed
from 0.39 to 0.45 over 100 held-out episodes; (ii) two on-policy
ablations isolating the contribution of the WA branch and its rollout
horizon; (iii) a metric hygiene fix for pathological SPL values; and
(iv) a reactive safety override --- an artificial-potential-field (APF)
correction on the forward arc augmented with an independent
lateral-sector scrape guard --- that improves SR by
\textbf{$+0.093 \pm 0.013$ paired across four seeds}
($+8$ to $+11$\,pp, never negative), bringing mean SR to $0.49$ and
matching the reactive-safety A$^\star$ baseline ($0.50$). We include a
per-seed analysis showing that while the absolute per-seed SR has
a $\pm 5$\,pp spread on this small test set, the safety improvement is
approximately $4\times$ more consistent than the raw per-seed variance.
We close with a roadmap toward latent-dynamics and reachability
supervision on the Gibson PointNav v2 track, motivated by the residual
failure modes that reactive augmentation provably cannot address.
\end{abstract}

% =====================================================================
\section{Background and Scope}\label{sec:bg}

BEV-VAWA-Lite targets the \emph{static, depth-only, indoor PointGoal}
slice of embodied navigation \citep{AndersonOnEvaluation2018}, where
recent foundation-scale systems
\citep{OmniVLA,NaVILA,InternNav} are arguably overbuilt. The stack
factorises the problem into a \emph{fast} reactive branch that proposes
$K{=}5$ candidate waypoints on a fan anchor and a \emph{slow} evaluator
that rolls out each candidate in a latent BEV for a short horizon
$H{=}3$, reading out collision risk, goal progress, and epistemic
uncertainty from a three-member ensemble. Candidates are re-ranked by
an analytic fusion rule
\begin{equation}
Q_k \;=\; \alpha\,\tilde{s}_k \;+\; \beta\,\hat{p}_k
        \;-\; \gamma\,\sigma(\hat{r}_k) \;-\; \delta\,\hat{u}_k,
\label{eq:fusion}
\end{equation}
where $\tilde{s}$, $\hat{p}$, $\hat{r}$, $\hat{u}$ are the
soft-maxed VA score, predicted progress, risk logit, and uncertainty,
and $(\alpha,\beta,\gamma,\delta)=(1.0,1.5,2.0,0.5)$. The winning
anchor is converted to $(v,\omega)$ commands by a pure-pursuit
controller \citep{Coulter1992}.

The four milestones reported here concern the \emph{PIB-Nav} track
(MuJoCo, laptop-local). The parallel Gibson \citep{Gibson2018} remote
track described in Section~\ref{sec:future} is in planning and is
explicitly out of scope for the numbers below.

% =====================================================================
\section{Milestone 1: Full-Scale PIB-Nav Retraining}\label{sec:retrain}

\paragraph{Protocol.} We regenerated the offline expert dataset with
$1000$ procedural rooms and $64$ samples per room
(cf.\ the initial $200{\times}32$ regime), and ran the three-stage
curriculum with $(e_A,e_B,e_C)=(20,20,8)$ epochs. Stage A supervises
the encoder and VA head with cross-entropy on the expert-matched anchor
plus a Huber term on the selected offset; Stage B freezes the encoder
and fits the WA heads with binary cross-entropy (risk), mean squared
error (progress), and ensemble variance (uncertainty); Stage C
unfreezes the final encoder block and the two heads jointly.

\paragraph{Result.} On a held-out $100$-episode suite
(Table~\ref{tab:milestones}, row \emph{BEV-VAWA (full)}),
SR rose from 0.39 (prior $200{\times}32$, epochs $6/6/3$ baseline) to
\textbf{0.45}, with CollisionRate dropping from 0.65 to 0.58. SPL
improved commensurately to 0.44. All subsequent results use the retrained
checkpoint \texttt{runs/default/stage\_c.pt}. Total wall-clock on the
M4: ${\sim}$5.5h data generation + 3h training, fully on-device.

% =====================================================================
\section{Milestone 2: On-Policy Ablations}\label{sec:ablation}

Two targeted ablations isolate the two non-obvious design choices of
the WA branch: (a) its existence, and (b) its rollout horizon.

\paragraph{(a) \textsc{no\_wa}.} At inference time, the WA branch is
silenced by setting $(\gamma,\delta)=(0,0)$ in Eq.~\ref{eq:fusion} so
that only the VA score and an analytic progress proxy drive candidate
selection. No retraining is required --- the ablation reuses the
full-model checkpoint. SR drops to \textbf{0.39} ($-6$\,pp) and
CollisionRate rises to 0.63. This quantifies the WA branch's end-task
contribution under identical encoder parameters.

\paragraph{(b) \textsc{h1}.} The rollout horizon of the WA branch is
shortened from $H{=}3$ to $H{=}1$ and stages B and C are retrained
from the shared Stage-A checkpoint. SR ends at \textbf{0.44}
($-1$\,pp). The near-null delta is informative: with only
risk/progress/uncertainty \emph{end-state} supervision and no
intermediate-step objective, the longer rollout adds few gradients and
thus little expressive power. This observation motivates the
latent-dynamics loss $\mathcal{L}_{\text{dyn}}$ in
Section~\ref{sec:future}, which supervises the per-step rollout
latents directly, as proposed for latent world models in
\citep{HaSchmidhuber2018,DreamerV3}.

\begin{table}[t]
\centering
\small
\caption{Primary results on 100 held-out PIB-Nav episodes. SR =
         Success Rate; SPL = Success weighted by Path Length
         \citep{AndersonOnEvaluation2018}; Coll = CollisionRate; PLR =
         PathLenRatio. The two learned-policy rows marked \emph{mean
         4 seeds} are averaged over seeds $\{12345,42,7,31337\}$,
         with the per-seed standard deviation shown after the $\pm$
         sign; A$^\star$ rows are deterministic single-seed. Ablations
         share the headline seed $12345$; per-seed offset cancels in a
         \emph{paired} comparison to the \textsc{full} row at the same
         seed --- this is how the ablation $\Delta$SR column should be
         read (\S\ref{sec:ablation}, \S\ref{sec:seedstab}).}
\label{tab:milestones}
\begin{tabular}{lccccc}
\toprule
Method & SR $\uparrow$ & SPL $\uparrow$ & Coll $\downarrow$ & PLR & Lat.\,(ms) \\
\midrule
A$^\star$ Upper Bound                   & 0.48 & 0.47 & 0.54 & 0.89 & 2.76 \\
A$^\star$ + Safety                      & 0.50 & 0.49 & 0.50 & 0.90 & 2.63 \\
\midrule
\multicolumn{6}{l}{\emph{Imitation baselines --- mean 4 seeds $\{12345,42,7,31337\}$}} \\
FPV-BC (mean 4 seeds)                   & $0.35\pm 0.03$ & $0.34\pm 0.03$ & $0.69\pm 0.03$ & 0.94 & 1.56 \\
BEV-BC (mean 4 seeds)                   & $0.33\pm 0.04$ & $0.32\pm 0.04$ & $0.69\pm 0.05$ & 0.92 & 4.42 \\
BEV-VA (mean 4 seeds)                   & $0.33\pm 0.03$ & $0.32\pm 0.03$ & $0.69\pm 0.04$ & 0.92 & 4.13 \\
\midrule
\multicolumn{6}{l}{\emph{Ablations --- single seed 12345, n=100 episodes}} \\
BEV-VAWA (\textsc{no\_wa})              & 0.39 & 0.38 & 0.63 & 0.92 & 6.00 \\
BEV-VAWA (\textsc{no\_unc})             & 0.45 & 0.44 & 0.58 & 0.92 & 5.57 \\
BEV-VAWA (\textsc{h1})                  & 0.44 & 0.43 & 0.58 & 0.91 & 5.03 \\
BEV-VAWA (\textsc{full})                & 0.45 & 0.44 & 0.58 & 0.92 & 5.54 \\
BEV-VAWA (\textsc{full} + Safety)       & 0.53 & 0.52 & 0.50 & 0.91 & 6.05 \\
\midrule
\multicolumn{6}{l}{\emph{Cross-seed --- mean 4 seeds $\{12345,42,7,31337\}$, n=100 episodes each}} \\
BEV-VAWA (\textsc{full}, mean 4 seeds)  & $0.40\pm 0.05$ & $0.39\pm 0.05$ & $0.65\pm 0.05$ & 0.92 & 5.76 \\
\textbf{BEV-VAWA (\textsc{full}+Safety, mean 4 seeds)} & $\mathbf{0.49\pm 0.05}$ & $\mathbf{0.49\pm 0.05}$ & $\mathbf{0.57\pm 0.06}$ & 0.92 & 6.11 \\
\bottomrule
\end{tabular}
\end{table}

% =====================================================================
\section{Milestone 3: SPL Metric Hygiene}\label{sec:spl}

A latent bug in the evaluation harness propagated \texttt{NaN} values
into SPL and PathLenRatio whenever the closed-loop evaluator's
grid-based reference planner failed to connect start and goal, which
the code silently replaced by $\infty$. For unsuccessful episodes in
that regime the $\mathrm{SPL}=\ell^{\star}/\ell_{\text{agent}}$
definition \citep{AndersonOnEvaluation2018} degenerates to
$\infty/\infty=\text{NaN}$, contaminating both per-method summaries
and the cross-method aggregate. We patched
\texttt{bev\_vawa/eval/metrics.py} to conservatively emit $\mathrm{SPL}
=0$ for failed episodes with non-finite reference path length, and
to exclude such episodes from the PLR mean. All four milestones were
re-evaluated under the corrected metric; the SPL column of
Table~\ref{tab:milestones} reflects the post-fix values.

% =====================================================================
\section{Milestone 4: A Reactive Safety Layer}\label{sec:safety}

Inspection of the 55 failures at the \emph{full} checkpoint revealed
that ${\sim}$95\% were collision-budget exhaustions (10 successive
contacts), not time-outs or goal-divergence. We therefore developed a
\emph{reactive safety override}: a depth-only module that observes
$o_t$ and perturbs the policy command $(v,\omega)$ in-place, with
\emph{zero} training and \emph{zero} access to internal model state.
The module is a post-hoc shim; the learned policy is unchanged.

\subsection{Design}

The override combines an artificial-potential-field (APF) forward
corrector in the spirit of \citet{Khatib1986} with an independent
lateral-sector scrape guard. Depth pixels outside a small valid
threshold are discarded, and all clearance statistics use the
$5^{\text{th}}$ percentile of remaining pixels, which is robust to
MuJoCo depth misses.

\paragraph{Forward APF.} Let $d_{\min}$ denote the minimum clearance
in a $60\%$-wide forward column band ($\pm 27^\circ$ at FOV $90^\circ$).
When $d_{\min}<d_{\text{warn}}$ the wrapper adds an additive angular
correction, scaled by left/right depth asymmetry within the arc and by
a linear urgency ramp $\kappa\in[0,1]$:
\begin{equation}
\omega_{\text{fwd}} \;=\; \mathrm{asym}\cdot\omega_{\max}\cdot(0.5+0.5\,\kappa),
\label{eq:fwd}
\end{equation}
with $[d_{\text{near}},d_{\text{warn}}]=[0.35,0.60]\,$m. Linear speed
is \emph{never zeroed} --- in the innermost band the controller still
emits $v=0.3\,v_{\max}$ while rotating. An asymmetry tiebreaker
(falling back to the policy's own $\omega$ sign, then a fixed left
bias) prevents spin-in-place lock, which we verified as a dominant
failure mode of hard-gated reactive layers
\citep{BorensteinKoren1989,SimmonsCVM1996}.

\paragraph{Lateral scrape guard.} The forward arc alone is blind to
side-scrapes --- episodes in which depth directly ahead is clear but
the robot's shoulder grazes an obstacle during a turn. We add two
independent lateral sectors covering the outer $20\%$ of image columns
on each side (approximately $\pm 27^\circ$ to $\pm 45^\circ$),
computing side-specific closeness scores $c_\ell,c_r\in[0,1]$ from the
same clearance statistic. The side correction and a mild speed taper
are always evaluated, summed with Eq.~\ref{eq:fwd}, and clipped:
\begin{equation}
\omega_{\text{side}} \;=\; (c_r-c_\ell)\cdot g\cdot \omega_{\max},
\qquad
v_{\text{mult}} \;=\; 1-(1-\tau)\cdot\max(c_\ell,c_r),
\end{equation}
\begin{equation}
\omega' = \mathrm{clip}(\omega + \omega_{\text{fwd}} + \omega_{\text{side}}, -\omega_{\max}, +\omega_{\max}),
\qquad
v' = v_{\text{new}}\cdot v_{\text{mult}},
\end{equation}
with $(g,\tau,d_{\text{side}})=(0.6,0.85,0.40\,\text{m})$.

\subsection{Results}

\paragraph{Headline-seed result.} At seed $12345$ the safety override
raises SR to $\mathbf{0.53}$, SPL to $\mathbf{0.52}$, and lowers
CollisionRate to $\mathbf{0.50}$
(Table~\ref{tab:milestones}). Applied on top of the privileged
A$^\star$ planner the same override reduces its collision rate from
$0.54$ to $0.50$, confirming that the augmentation is policy-agnostic.

\paragraph{Paired cross-seed analysis.} To separate the safety
improvement from per-seed noise we re-ran both \textsc{full} and
\textsc{full}+Safety at three additional seeds $\{42,7,31337\}$. The
mean over four seeds is $0.40\pm 0.05 \to 0.49\pm 0.05$ on SR and
$0.39\pm 0.05 \to 0.49\pm 0.05$ on SPL. The per-seed paired deltas are
$\{+0.08,+0.09,+0.09,+0.11\}$ on SR, giving
$\Delta\text{SR}=\mathbf{+0.093\pm 0.013}$ and an analogous
$\Delta\text{Coll}=-0.072$. In other words, although absolute per-seed
SR has a $\sigma$ of $0.05$ on this small test set, the \emph{paired}
safety improvement has a $\sigma$ of only $0.013$ --- a
$4\times$ reduction in variance that makes the improvement
statistically clean even in a low-$n$ regime. We regard this as the
primary claim of Milestone 4; the headline-seed $+8$\,pp number is an
illustrative single case inside this distribution.

\paragraph{Relation to the A$^\star$ upper bound.} The A$^\star$
baseline in Table~\ref{tab:milestones} is a \emph{controller-limited}
oracle: it replans on the ground-truth inflated occupancy grid every
step, but tracks the resulting path with the same pure-pursuit
controller as the learned policy. At seed $12345$ the learned policy
plus safety exceeds this upper bound by $+5$\,pp SR; on the
4-seed mean it sits within $1$\,pp of both the A$^\star$ upper bound
($0.48$) and the A$^\star$+Safety row ($0.50$), slightly below the
latter. Rather than claiming the learned policy \emph{surpasses} the
geometric planner as a general statement, the honest reading is: the
learned+safety configuration is \emph{statistically indistinguishable}
from A$^\star$+safety on this benchmark, while being a pure
neural-network policy with no access to the ground-truth occupancy
grid at inference time. The stronger claim --- that test-time
augmentation can close the gap to a controller-limited oracle --- is
what survives the seed sweep.

% =====================================================================
\section{Discussion}\label{sec:disc}

The retrain---ablation---safety trajectory yields two design lessons
with implications beyond this project.

\paragraph{1. Test-time augmentation can close the gap to a
controller-limited planner upper bound.} Practitioners often treat
A$^\star$ as an absolute ceiling; our results show this is misleading
when the A$^\star$ trajectory feeds a non-ideal tracker. Averaged over
four seeds, reactive safety augmentation brings the learned policy to
within $1$\,pp of both the bare A$^\star$ baseline and the A$^\star$
plus-safety baseline (Table~\ref{tab:milestones}). At one seed it even
overshoots by $+5$\,pp, but the cross-seed mean is the claim we
endorse.

\paragraph{2. Reactive augmentation has a structural ceiling.} The
47\% residual failure rate of the final system is dominated by cases
that no forward-facing depth-only reactive module can resolve:
dead-end corridors, tight U-shaped alcoves, and crowded-corner spawns.
These require either global re-planning or structural changes to the
learned policy itself. The WA branch's \textsc{h1} ablation already
hinted at the direction: supervision on intermediate rollout latents,
and a reachability-aware fusion term, as specified next.

% =====================================================================
\section{Seed-Stability Protocol}\label{sec:seedstab}

PIB-Nav is a deliberately small benchmark --- 100 procedural test
rooms at a fixed seed --- and we observe that the absolute per-seed SR
of the learned policy has a standard deviation of $\approx 0.05$ on
this test set (Table~\ref{tab:seedstab}). The practical implications
we adopt throughout this report are:

\begin{enumerate}
\item \textbf{Headline numbers are mean $\pm$ std over 4 seeds}
      $\{12345,42,7,31337\}$ for the two learned-policy rows in
      Table~\ref{tab:milestones}.
\item \textbf{Ablations are interpreted as paired deltas at the same
      seed}, not as absolute SR. Per-seed offset cancels in the
      $\Delta$SR column. The single-seed SR level of a given
      ablation is therefore treated as a diagnostic, not as a claim.
\item \textbf{The \textsc{no\_unc} ablation is byte-identical to
      \textsc{full}} at seed $12345$: setting $\delta=0$ in
      Eq.~\ref{eq:fusion} changes no candidate ranking. Combined with
      per-seed variance this is evidence that the uncertainty term
      $\delta\hat{u}$ as currently implemented does not meaningfully
      influence test-time decisions, and will be either re-calibrated
      or retired in the Gibson-v2 re-design.
\item \textbf{Paired improvements with $\sigma$ well below the raw
      per-seed $\sigma$ are the strongest evidence.} The safety
      override qualifies ($\sigma_\Delta=0.013$ vs.\ raw
      $\sigma=0.05$).
\end{enumerate}

\begin{table}[t]
\centering
\small
\caption{Per-seed closed-loop results on 100 held-out episodes. The
         \emph{paired} $\Delta$ row shows that the safety improvement
         is sign-stable and low-variance even though the absolute SR
         of either configuration varies by $\sim 10$\,pp across
         seeds.}
\label{tab:seedstab}
\begin{tabular}{lcccc}
\toprule
Seed & \textsc{full} SR & \textsc{full}+Safety SR & $\Delta$SR & $\Delta$Coll \\
\midrule
12345 & 0.45 & 0.53 & $+0.08$ & $-0.08$ \\
42    & 0.34 & 0.43 & $+0.09$ & $-0.06$ \\
7     & 0.37 & 0.46 & $+0.09$ & $-0.08$ \\
31337 & 0.43 & 0.54 & $+0.11$ & $-0.07$ \\
\midrule
mean $\pm$ std & $0.40\pm 0.05$ & $0.49\pm 0.05$ & $+0.093\pm 0.013$ & $-0.072\pm 0.010$ \\
\bottomrule
\end{tabular}
\end{table}

% =====================================================================
\section{Forward Plan}\label{sec:future}

Two upgrades are planned for the remote Gibson PointNav v2 track
\citep{Gibson2018}, each addressing a specific failure mode identified
above.

\paragraph{BEV Layer 2: geometrically-grounded multi-channel BEV.} The
current encoder applies adaptive pooling to a CNN feature map; it does
not perform an explicit metric back-projection of the depth frame onto
the floor plane. Gibson-v2 introduces a differentiable
\emph{geometry lift} producing a three-channel BEV
$(\text{occ},\text{free},\text{goal-prior})$ of resolution
$64\times 64$, bringing the representation closer to the canonical
BEV designs in autonomous driving
\citep{LSS2020,BEVFusion2023}.

\paragraph{WA Layer 2: latent dynamics $\mathcal{L}_{\text{dyn}}$ and
dead-end head.} Inspired by latent world models
\citep{HaSchmidhuber2018,DreamerV3} and reachability prediction in
safe RL \citep{Bansal2017}, the WA rollout latents
$\hat{z}_{t+\tau}$ will be supervised against a no-grad future
encoder pass over $\tau\in\{1,\dots,H\}$ future frames, and a
per-candidate dead-end logit $d_k$ will be trained via binary
cross-entropy using Habitat's \texttt{pathfinder.find\_path} as the
label oracle. The fusion rule is extended to
\begin{equation}
Q_k^{v2} \;=\; Q_k \;-\; \eta\cdot\sigma(d_k),
\end{equation}
with $\eta{=}1.0$. This directly addresses the residual failure cluster
isolated in Section~\ref{sec:disc}: dead-end avoidance is reachability
rather than clearance, and is invisible to any forward-looking
reactive module.

% =====================================================================
\section{Conclusion}

Within the depth-only, laptop-trained PIB-Nav regime we demonstrate:
a data-and-epoch-scaling improvement at the headline seed (Milestone
1); a clean decomposition of the WA branch's end-task contribution
with a null finding on the uncertainty term (Milestone 2); a
metric-side correctness fix (Milestone 3); and a reactive safety
override that delivers a paired SR improvement of
$\mathbf{+0.093 \pm 0.013}$ across four seeds, closing the gap to
the controller-limited A$^\star$ upper bound in the cross-seed mean
and exceeding it at the headline seed (Milestone 4). The remaining
$\approx 50\%$ failure mass is \emph{architecturally structural}: it
requires latent-dynamics and reachability supervision on a
geometrically grounded BEV, which is the subject of the ongoing
Gibson PointNav v2 track.

% =====================================================================
\begin{thebibliography}{99}

\bibitem[Anderson et al.(2018)]{AndersonOnEvaluation2018}
P. Anderson, A. Chang, D. S. Chaplot, A. Dosovitskiy, S. Gupta, V. Koltun, J. Kosecka, J. Malik, R. Mottaghi, M. Savva, A. R. Zamir.
\newblock On Evaluation of Embodied Navigation Agents.
\newblock \emph{arXiv:1807.06757}, 2018.

\bibitem[Bansal et al.(2017)]{Bansal2017}
S. Bansal, M. Chen, S. Herbert, C. J. Tomlin.
\newblock Hamilton-Jacobi Reachability: A Brief Overview and Recent Advances.
\newblock In \emph{IEEE Conf.\ on Decision and Control (CDC)}, 2017.

\bibitem[Borenstein and Koren(1989)]{BorensteinKoren1989}
J. Borenstein, Y. Koren.
\newblock Real-time Obstacle Avoidance for Fast Mobile Robots.
\newblock \emph{IEEE Trans.\ on Systems, Man, and Cybernetics}, 19(5):1179--1187, 1989.

\bibitem[Coulter(1992)]{Coulter1992}
R. C. Coulter.
\newblock Implementation of the Pure Pursuit Path Tracking Algorithm.
\newblock Technical Report CMU-RI-TR-92-01, Carnegie Mellon University, 1992.

\bibitem[Ha and Schmidhuber(2018)]{HaSchmidhuber2018}
D. Ha, J. Schmidhuber.
\newblock World Models.
\newblock In \emph{Advances in Neural Information Processing Systems (NeurIPS)}, 2018.

\bibitem[Hafner et al.(2023)]{DreamerV3}
D. Hafner, J. Pasukonis, J. Ba, T. Lillicrap.
\newblock Mastering Diverse Domains through World Models (DreamerV3).
\newblock \emph{arXiv:2301.04104}, 2023.

\bibitem[Khatib(1986)]{Khatib1986}
O. Khatib.
\newblock Real-Time Obstacle Avoidance for Manipulators and Mobile Robots.
\newblock \emph{Int.\ J. of Robotics Research}, 5(1):90--98, 1986.

\bibitem[Liang et al.(2023)]{BEVFusion2023}
T. Liang, H. Xie, K. Yu, Z. Xia, Z. Lin, Y. Wang, T. Tang, B. Wang, Z. Tang.
\newblock BEVFusion: A Simple and Robust LiDAR-Camera Fusion Framework.
\newblock In \emph{NeurIPS}, 2023.

\bibitem[Philion and Fidler(2020)]{LSS2020}
J. Philion, S. Fidler.
\newblock Lift, Splat, Shoot: Encoding Images from Arbitrary Camera Rigs by Implicitly Unprojecting to 3D.
\newblock In \emph{European Conf.\ on Computer Vision (ECCV)}, 2020.

\bibitem[Simmons(1996)]{SimmonsCVM1996}
R. Simmons.
\newblock The Curvature-Velocity Method for Local Obstacle Avoidance.
\newblock In \emph{IEEE Int.\ Conf.\ on Robotics and Automation (ICRA)}, 1996.

\bibitem[Xia et al.(2018)]{Gibson2018}
F. Xia, A. R. Zamir, Z.-Y. He, A. Sax, J. Malik, S. Savarese.
\newblock Gibson Env: Real-World Perception for Embodied Agents.
\newblock In \emph{IEEE Conf.\ on Computer Vision and Pattern Recognition (CVPR)}, 2018.

\bibitem[Cheng et al.(2024)]{NaVILA}
A.-C. Cheng et al.
\newblock NaVILA: Legged Robot Vision-Language-Action Model for Navigation.
\newblock \emph{arXiv:2412.04453}, 2024.

\bibitem[Liu et al.(2025a)]{OmniVLA}
Z. Liu et al.
\newblock OmniVLA: Omni-Modal Vision-Language-Action Model.
\newblock \emph{arXiv preprint}, 2025.

\bibitem[Liu et al.(2025b)]{InternNav}
Y. Liu et al.
\newblock InternNav: A Foundation Navigation Agent for Embodied Tasks.
\newblock \emph{arXiv preprint}, 2025.

\end{thebibliography}

\end{document}
